\newpage
\section{Cronograma de atividades do TCC02}
\label{sc:cronograma}

\begin{longtable}{p{\textwidth}}
\label{tab:fases}\\
\toprule%
\myrowcolour%
\bfseries F.1: Revisão bibliográfica \\
\midrule
Aprofundamento e atualização da revisão iniciada no TCC01, fixando o referencial sobre mapeamento objeto-relacional, o padrão \emph{Data Mapper} e a metaprogramação em TypeScript que fundamenta a implementação.
\\
\midrule
\myrowcolour%
\bfseries Atividades planejadas para a fase 1 \\
\midrule
A.1.1 Levantar e atualizar a literatura sobre ORMs, \emph{Data Mapper}, \emph{decorators} e programação em nível de tipos. \\
A.1.2 Consolidar o comparativo do OOR com os ORMs concorrentes (TypeORM, Prisma, Drizzle, Sequelize). \\
\midrule
\myrowcolour%
\bfseries Resultados esperados para a fase 1 \\
\midrule
Base bibliográfica consolidada e comparativo de referência que orienta as decisões de projeto das fases seguintes. \\
\toprule%
\myrowcolour%
\bfseries F.2: \emph{Decorators} para o ORM OOR com \emph{Data Mapper} \\
\midrule
Implementação da camada de definição de entidades por \emph{decorators} e da camada de persistência segundo o padrão \emph{Data Mapper}, atendendo aos objetivos específicos de declaração do mapeamento, tipagem em compilação e persistência por repositório.
\\
\midrule
\myrowcolour%
\bfseries Atividades planejadas para a fase 2 \\
\midrule
A.2.1 Implementar os \emph{decorators} (\texttt{@Entity}, \texttt{@Column}, \texttt{@PrimaryColumn}, \texttt{@ToOne}, \texttt{@Nullable}) e o \texttt{MetadataStorage} por \texttt{Database}, sem \texttt{reflect-metadata}. \\
A.2.2 Implementar o \texttt{Repository<T>} com as operações de persistência sobre cada entidade. \\
A.2.3 Aplicar programação em nível de tipos (\emph{brand} \texttt{PrimaryKey<V>}, ausência de \texttt{any}) e dar suporte à herança de tabela única (STI). \\
\midrule
\myrowcolour%
\bfseries Resultados esperados para a fase 2 \\
\midrule
Entidades declaráveis por \emph{decorators} que verifica em tempo de compilação possíveis erros de mapeamento e suporte a STI.
Entidades persistíveis via repositório. \\
\toprule%
\myrowcolour%
\bfseries F.3: \emph{Query builder} para o ORM OOR \\
\midrule
Construção do mecanismo de composição de consultas, que permite filtrar e recuperar entidades por encadeamento de operações --- sem SQL em texto --- e traduzi-las para SQL sempre parametrizado.
\\
\midrule
\myrowcolour%
\bfseries Atividades planejadas para a fase 3 \\
\midrule
A.3.1 Implementar o \texttt{QueryBuilder<T>} \emph{lazy}, com a interface de composição encadeável de consultas. \\
A.3.2 Traduzir as consultas para SQL parametrizado por meio do \emph{driver} \texttt{SQL} do Bun. \\
A.3.3 Garantir a verificação das consultas pelo sistema de tipos em compilação. \\
\midrule
\myrowcolour%
\bfseries Resultados esperados para a fase 3 \\
\midrule
Leituras expressas sem SQL em texto, verificadas em compilação e traduzidas para consultas parametrizadas, completando o núcleo funcional do ORM. \\
\toprule%
\myrowcolour%
\bfseries F.4: Teste de desempenho ou de usabilidade \\
\midrule
Avaliação experimental do OOR, contrastando-o com ORMs concorrentes em desempenho e/ou em usabilidade da API, tomando como referência metodológica o \emph{benchmark} revisado.
\\
\midrule
\myrowcolour%
\bfseries Atividades planejadas para a fase 4 \\
\midrule
A.4.1 Definir a metodologia, as métricas e os cenários de teste, bem como os ORMs de comparação. \\
A.4.2 Executar os testes e coletar as medições de forma reprodutível. \\
\midrule
\myrowcolour%
\bfseries Resultados esperados para a fase 4 \\
\midrule
Conjunto de dados experimentais coletados sob protocolo definido, comparáveis entre o OOR e os ORMs concorrentes. \\
\toprule%
\myrowcolour%
\bfseries F.5: Análise do teste de desempenho ou de usabilidade \\
\midrule
Tratamento, interpretação e discussão dos dados coletados na fase anterior, situando o OOR frente aos concorrentes e às hipóteses do trabalho.
\\
\midrule
\myrowcolour%
\bfseries Atividades planejadas para a fase 5 \\
\midrule
A.5.1 Tabular os dados e aplicar as medidas de resumo (estatística descritiva, gráficos). \\
A.5.2 Analisar e discutir os resultados em comparação com os ORMs concorrentes. \\
\midrule
\myrowcolour%
\bfseries Resultados esperados para a fase 5 \\
\midrule
Análise consolidada que evidencia os pontos fortes e as limitações do OOR e responde às questões de pesquisa. \\
\toprule%
\myrowcolour%
\bfseries F.6: Redação do artigo \\
\midrule
Redação do documento final do TCC02, integrando fundamentação, implementação e resultados, com revisão pelo orientador e preparação da defesa.
\\
\midrule
\myrowcolour%
\bfseries Atividades planejadas para a fase 6 \\
\midrule
A.6.1 Redigir o artigo, articulando a revisão, os métodos, a implementação e os resultados. \\
A.6.2 Revisar o texto com o orientador e incorporar os ajustes. \\
A.6.3 Preparar a apresentação e a defesa perante a banca. \\
\midrule
\myrowcolour%
\bfseries Resultados esperados para a fase 6 \\
\midrule
Artigo final concluído e revisado, pronto para submissão e defesa. \\
\bottomrule
\end{longtable}


\begin{ganttchart}[
y unit title=.6cm,
y unit chart=.6cm,
x unit = 2.8cm,
hgrid,
vgrid,
time slot format=isodate-yearmonth,
time slot unit=month,
inline,
bar height=0.7,
bar/.append style={fill=gray!50, inner sep=0pt}
]{2026-08}{2026-12}
\gantttitlecalendar{year, month} \\
  \ganttbar{F.1: Revisão bibliográfica}{2026-08}{2026-09} \\
  \ganttbar{F.2: \emph{Decorators}}{2026-08}{2026-10} \\
  \ganttbar{F.3: \emph{Query builder}}{2026-10}{2026-11} \\
  \ganttbar{F.4: Testes}{2026-11}{2026-12} \\
  \ganttbar{F.5: Análise dos testes}{2026-12}{2026-12} \\
  \ganttbar{F.6: Redação do artigo}{2026-11}{2026-12} \\
\end{ganttchart}